\chapter*{Conclusion Générale}
\markboth{Conclusion Générale}{}
\addcontentsline{toc}{chapter}{Conclusion Générale}

Ce projet de fin d'études avait pour objectif de doter la plateforme Netethic d'une couche
d'intelligence artificielle capable d'accompagner les établissements scolaires dans leur lutte
contre le harcèlement numérique. L'enjeu était de concevoir des modules adaptés non seulement
au langage naturel, mais aussi au contexte émotionnel et juridique propre à ce domaine.

Les quatre modules réalisés répondent chacun à un besoin précis : \textbf{Noa} assure une présence
permanente sur le site public, l'\textbf{assistant technique} accompagne le personnel dans la navigation et l'utilisation
de la plateforme, l'\textbf{assistant émotionnel} aide les etudiants à exprimer leurs émotions et à gérer leurs conflits, 
et le \textbf{classificateur de harcèlement},
ancré dans le droit français, offre une lecture structurée des contenus signalés selon quatre
stratégies complémentaires. L'ensemble forme un système cohérent, modulaire et conteneurisé,
dont les résultats valident la pertinence des choix effectués.

Ce travail ouvre néanmoins des perspectives ambitieuses. La \textbf{détection multimodale}
constitue un axe naturel : le harcèlement numérique emprunte aussi la forme d'images, de memes
et de messages vocaux, ce qui appelle une combinaison de vision par ordinateur et de traitement
du langage.Enfin, une
\textbf{analyse comportementale longitudinale} permettrait de passer d'une posture réactive à
une posture prédictive en alertant les équipes éducatives avant qu'une situation n'atteigne un
seuil critique, et faisant ainsi de Netethic non plus un outil de détection, mais de prévention.
